\section{Discussion}
\label{sec:discussion}

\subsection{Interpretation of the Main Result}

The experiments support the central hypothesis of the project: metadata
selection is not merely a prompt-compression mechanism. In the evaluated
setting, a smaller relation-aware and SQL-safe context produces better
held-out results than complete-catalog prompting while also reducing model
input cost.

The central interpretation is therefore qualitative as well as quantitative.
The language model benefits when semantic descriptions are preserved but
internal metadata identifiers and unrelated schema elements are prevented
from competing with the physical Spark objects needed for execution.


\subsection{Why the Complete Catalog Can Underperform}

At first sight, the complete catalog should appear advantageous because it
contains all information needed to answer every benchmark question. However,
the experiments show that information availability and information usefulness
are not equivalent.

The full catalog includes physical schema elements together with semantic
concepts, aliases, relationship identifiers, and rule identifiers. These
metadata objects are useful internally, but they are not all valid SQL table
or column names.

An LLM exposed to the complete catalog must therefore distinguish between
physical identifiers and descriptive metadata while simultaneously solving
the analytical task. The held-out failures show that the model can incorrectly
promote an internal semantic identifier into generated SQL.

The SQL-safe relation-aware renderer removes much of this ambiguity. It keeps
semantic meaning available through descriptions while exposing only physical
tables, columns, join conditions, and executable expressions as SQL objects.

The advantage of Configuration B therefore comes from both
\emph{selection} and \emph{representation}: less metadata is shown, and the
metadata that remains is presented in a form that more directly matches the
execution language.


\subsection{Retrieval Completeness and SQL Accuracy}

One of the clearest findings is the association between metadata recall and
downstream query correctness.

For held-out questions with perfect metadata recall, the relation-aware
configuration achieves 88.2\% result accuracy, while the validation-and-repair
configuration reaches 94.1\%.

For questions with imperfect metadata recall, accuracy drops to 33.3\% for
both configurations.

This difference is expected from the system architecture. The LLM can only
generate a correctly grounded cross-source query if its prompt contains the
physical datasets, columns, relationships, and semantic rules needed to
express that query.

When required metadata is missing, the model may invent a plausible schema
element or simplify the query so that an intended source is omitted. Such
errors cannot always be detected from syntax alone.

This suggests that metadata recall is a critical upstream metric for
metadata-aware Text-to-SQL systems. High retrieval precision is useful for
keeping prompts compact, but insufficient recall can directly remove
information required for executable query construction.


\subsection{Limits of Semantic Similarity}

The catalog-scale experiment highlights a limitation of dense semantic
retrieval.

At four sources, the relation-aware retriever achieves macro recall of 0.925.
After four domain-related hard-negative sources are introduced, macro recall
drops to 0.825.

The error analysis shows that the new failures are not random. They are
concentrated around geographic metadata, especially Taxi Zones relationships
and borough-related fields.

This result illustrates a general challenge in metadata retrieval:
semantically similar sources can compete for a small Top-\(k\) retrieval
budget even when only one source has the correct structural relationship
required for execution.

The deterministic relation-expansion stage mitigates this problem only after
the relevant dataset or concept has been selected. If dense and lexical
retrieval fail to provide enough evidence for the correct dataset in the
first place, later relationship expansion cannot reliably recover it.

The experiment therefore shows that semantic similarity and structural
correctness are complementary rather than interchangeable.


\subsection{Catalog Growth and Bounded Context}

A positive scalability result is that relation-aware context size remains
approximately bounded while the full catalog grows substantially.

The complete catalog grows from 9,498 to 16,651 characters between four and
sixteen sources. The retrieved context remains near 2,100 characters.

This behaviour is desirable because the per-query prompt does not need to
scale linearly with total metadata volume.

However, the corresponding increase in context reduction must not be
interpreted in isolation. At eight and sixteen sources, part of the greater
compression is associated with missing required metadata.

The correct interpretation is therefore a trade-off: relation-aware retrieval
successfully decouples prompt size from catalog size, but aggressive selection
can reduce recall when hard negatives compete with the correct metadata.

A production system would need to manage this trade-off explicitly, for
example by using adaptive retrieval depth, stronger structural reranking, or
confidence-aware expansion.


\subsection{Interpreting the 8-to-16 Source Plateau}

An interesting aspect of the catalog-scale experiment is that performance does
not continue to degrade when the catalog increases from eight to sixteen
sources.

Macro recall, precision, F1, perfect-recall count, and mean context size are
identical at these two scales.

This suggests that catalog size alone is not the main cause of the observed
failure. Instead, the first group of hard negatives contains metadata that is
particularly competitive with the intended sources.

The additional eight sources increase the total catalog substantially but do
not introduce new retrieval errors for the frozen benchmark.

This distinction matters because it prevents an overly broad conclusion such
as ``retrieval quality decreases as the number of sources increases.'' The
results support a more precise statement: retrieval quality can degrade when
new sources are semantically similar enough to displace structurally required
metadata.


\subsection{Effectiveness and Limits of Validation}

SQL validation provides a useful safety boundary between the language model
and Spark execution.

The validator can detect unknown physical columns, invalid relations,
syntactic problems, incompatible expressions, and unexpected output
structures before a query is executed.

This protects the execution layer from a class of common LLM generation
errors. It also provides actionable feedback that can be passed to the repair
stage.

However, validation cannot establish full semantic correctness.

A valid query may still use the wrong filter, choose the wrong dataset,
produce a different categorical convention, or express a related but
incorrect analytical computation. These queries may pass parsing and Spark
analysis while returning an incorrect benchmark result.

For this reason, the project treats result accuracy as the primary correctness
metric rather than validation rate.


\subsection{Effectiveness and Limits of One-Shot Repair}

The repair stage provides a measurable but limited improvement.

Three held-out queries triggered repair, and one repair produced a validated
and correct result. This raises overall result accuracy from 80\% to 85\%.

The successful case demonstrates the intended role of repair: the necessary
metadata was already available, the generated query contained a physical
schema error, and validator feedback gave the model enough information to
correct the query.

The failed repair cases reveal the opposite condition. When retrieval omitted
essential weather or relationship metadata, the repair model was given no new
source of structural knowledge. A validation error can identify that a column
does not exist, but it cannot reliably reconstruct an omitted multi-source
relationship.

One-shot repair is therefore best understood as a local correction mechanism,
not as a substitute for metadata retrieval.


\subsection{Prompt Reduction and Latency}

The relation-aware context reduces mean held-out prompt size by 69.66\%.

The controlled latency benchmark shows a similar reduction in prompt
evaluation cost: mean prompt-evaluation time decreases by 70.29\% under the
prefix-isolated protocol.

The similarity between these reductions is consistent with the expectation
that shorter prompts require less input processing by the language model.

The original sequential held-out timings, however, show why LLM latency must
be measured carefully. Repeated complete-catalog prompts share a large common
prefix and can benefit substantially from runtime cache behaviour.

Without the controlled benchmark, these cache effects could lead to the
incorrect conclusion that the larger full-catalog prompt is intrinsically
faster.

The separate prefix-isolated and natural-cache measurements therefore improve
the reliability of the efficiency analysis.


\subsection{Interpretation of the Spark Scaling Experiment}

The Spark data-volume experiment confirms that the complete analytical
pipeline can execute over progressively larger physical datasets, up to
2,963,713 curated Yellow Taxi records.

The temporal weather join shows the clearest runtime growth, increasing from
0.274 seconds at 100,000 rows to 1.350 seconds at full scale.

The simpler workloads show weaker monotonic behaviour because the experimental
datasets remain small enough for fixed execution costs and local caching
effects to be significant.

The results therefore demonstrate feasibility and relative behaviour within
the selected environment, but they should not be used to claim general Spark
scalability properties.

A distributed cluster experiment with substantially larger physical data
would be required to evaluate network shuffle, executor scaling, partitioning
strategy, and cluster-level resource utilisation.


\subsection{Big Data Perspective}

The project addresses the Big Data dimensions of Volume and Variety in
different parts of the system.

Volume is handled by Spark and evaluated by increasing the number of taxi
records processed by the physical execution layer.

Variety is handled primarily by the metadata architecture. Different physical
schemas, naming conventions, source roles, geographic joins, and temporal
granularities are represented explicitly in the catalog.

The experiments show that these two dimensions create different engineering
problems. Increasing row volume affects execution cost, whereas increasing
metadata variety affects grounding ambiguity and prompt construction.

Separating catalog scale from data scale is therefore useful when analysing
data-centric generative-AI systems: a system can scale well in physical query
execution while still failing because metadata grounding becomes ambiguous.


\subsection{Threats to Validity}

Several factors limit the generality of the experimental results.

First, the held-out benchmark contains 20 questions derived from ten analytical
intents. Although the natural-language paraphrases were frozen before final
evaluation, the benchmark remains small compared with large public
Text-to-SQL datasets.

Second, all core data comes from a single application domain: New York City
mobility and weather. The architecture may behave differently with metadata
from unrelated enterprise domains.

Third, the local language model is relatively small and executed without a
GPU. Larger models may have different sensitivity to metadata noise and may
require different retrieval depths.

Fourth, the catalog-scale distractors are synthetic metadata-only sources.
They are deliberately designed as hard negatives, but they do not reproduce
all characteristics of a large production catalog.

Fifth, the physical Spark evaluation is local rather than distributed and
reaches approximately three million rows. It therefore demonstrates prototype
behaviour rather than cluster-scale performance.

Finally, result correctness depends on the benchmark's expected-result and
comparison procedure. Some analytically equivalent outputs may differ in
labels or formatting, as observed in one cross-taxi comparison.


\subsection{Possible Improvements}

Several extensions could address the observed limitations without changing the
core design.

Retrieval could use an adaptive Top-\(k\) rather than a fixed value of five,
increasing metadata coverage when retrieval confidence is low.

A structural reranking stage could explicitly favour datasets connected by
catalog relationships required by the candidate analytical intent.

Retrieval confidence could also be exposed to the downstream pipeline. When
required metadata is uncertain, the system could request additional metadata
before generation rather than relying on repair after SQL failure.

The metadata catalog could be expanded with richer type constraints,
cardinality information, data-quality statistics, and source-level
provenance.

For larger deployments, the local SQLite and embedded Qdrant components could
be replaced by distributed metadata and vector services while preserving the
same logical interfaces.

Finally, repair could be extended beyond one attempt, although any iterative
strategy would require careful control to avoid uncontrolled agent loops and
to preserve reproducibility.


\subsection{Overall Assessment}

The prototype demonstrates that a relatively small local language model can
perform useful cross-source Spark SQL generation when it is supported by
explicit metadata grounding.

The key result is not that retrieval always succeeds. The catalog-scale
experiment shows clear cases in which hard negatives reduce recall.

Rather, the project demonstrates that the way metadata is selected and
rendered has a measurable effect on correctness, prompt size, and inference
cost.

The architecture therefore supports a data-centric interpretation of
Text-to-SQL: improving the representation and delivery of metadata can be as
important as changing the language model itself.
